\scanpage{90}
\mspage{- 175 -}
\hfill IX

\[
0\longrightarrow R^f/R^t\longrightarrow E/E^t\longrightarrow B\longrightarrow 0\quad . \tag*{(13.3.15)}
\]

On notera que
\[
R^f/R^t=(R_{\unsure{o|0}}/R_o^t)_S \tag*{(13.3.16)}
\]
est un groupe analytique constant sur $S$. Il résulte donc de 13.1.5 que $B$ est un groupe abéloide sur $S$ au voisinage de $s$ si et seulement si sa fibre en $s$
\[
B_o=A_o/T_o
\]
est un groupe analytique abéloide, de sorte que $A_o$ apparaît alors comme une extension d'un tel groupe $B_o$ par un tore $T_o$. \uline{Toute façon d'écrire} $A_o$ \uline{sous forme d'une extension d'un groupe abéloide par un tore complexe} (*) \uline{donne donc naissance sur} $A$ , \uline{au voisinage de} $s$, \uline{à une unique structure d'extension} (13.3.13) \uline{d'un groupe abéloide} $B$ \uline{par un tore} $T$, \uline{induisant la structure d'extension donnée en} $s$.

13.4. \uline{Interprétation en termes de} "\uline{structures de Hodge mixtes}" \uline{de DELIGNE}.

Reprenons d'abord le dictionnaire 13.1.4, que nous allons reformuler d'une façon légèrement différente,mieux adaptée à l'introduction des "structures de Hodge" pertinentes à la situation. Pour ceci, soit $\mathbb R$ le faisceau des sections de $R$ sur $S$, et introduisons le $\underline{O}_S$-Module
\[
\mathscr E(A)=\mathscr E=\underline R\otimes_{\mathbb Z}\underline O_S\quad , \tag*{(13.4.1)}
\]
qui est localement libre de type fini lorsque $R$ est localement constant. (NB Le faisceau $\mathscr E$ n'est pas nécessairement cohérent si $R$ n'est pas localement constant.) L'homomorphisme (13.1.3) $\underline R\longrightarrow\underline E$ définit alors

% Faint, partly covered numerals at lower strip. Attribution uncertain.
\margin{\unsure{493|1193} \illegible{1}}
% printed label at foot
493
% archive sheet number: 89
